\section{Burst time and burst size}
\label{sec:burst-time-size}

The two random quantities that a bursting cell presents to the outside world
are the time of its rupture and the size of the rupture. Both are now within
reach of closed forms.

\subsection{Burst time}
\label{sec:burst-time}

The no-burst probability is $I(t)=\Prob{\tau>t}$, so the density of the
burst time is
\begin{equation}
\varphi(t):=-I'(t)=\delta J(t),
\label{eq:phi}
\end{equation}
using $I'=-\delta J$ from \S\ref{sec:moments-recap}. This density is
\emph{defective}: since $I(t)\to b$,
\begin{equation}
\int_0^\infty\varphi(t)\,\dd t=1-b,
\end{equation}
with the remaining mass $b$ sitting at $\tau=\infty$ --- the realisations
that clear internally without ever bursting. Conditioned on eventual burst,
the density is $\varphi(t)/(1-b)=\delta J(t)/(1-b)$. Where earlier work
referred to ``the known pdf of rupture times $\phi(t)$'', it is
\eqref{eq:phi} that was meant, available all along from $I'=-\delta J$.

\figureflag{F4a.4}{Burst-time density}{%
\textbf{Purpose:} display the defective burst-time density, making the mass
$b$ at $\tau=\infty$ and the conditional density explicit in one picture.
\textbf{Type:} Python analytic plot. \textbf{Content:} two parameter sets,
$(\lambda,\mu,\delta)=(1,0.2,0.05)$ and $(1,0,0.1)$, on $t\in[0,15]$: the
defective density $\varphi(t)=\delta J(t)$ from \eqref{eq:phi} and
\eqref{eq:J} (solid curves); for the $\mu>0$ set, shade and annotate the
missing mass $b$ at $\infty$ (e.g.\ an annotated bar or hatched region
labelled ``mass $b$ at $\tau=\infty$''); overlay the conditional density
$\delta J(t)/(1-b)$ as a dashed curve for the $\mu>0$ set. \textbf{Axes:}
$t$ (horizontal), density (vertical); legend entries $\varphi(t)$ and
$\varphi(t)/(1-b)$ per parameter set. \textbf{Style:} unified Python style
--- default matplotlib with \texttt{tab:} colours, grid alpha 0.25,
\texttt{tight\_layout}, PDF output, fonts matching the chapter.
\textbf{Data source:} analytic formulas \eqref{eq:phi}, \eqref{eq:J}; no
simulation. \textbf{Placement:} \Secref{sec:burst-time}, directly after the
defect discussion.}

\subsection{Burst size}
\label{sec:burst-size}

The burst size $\Kb$ is the value of $\xt$ immediately before rupture. For
the burst to occur in state $k$, the process must be in state $k$ and then
rupture from there; rupturing from state $k$ happens at rate $\delta k$.
Hence
\begin{equation}
\pr{\Kb = k} = \delta k \int_0^\infty p_k(t)\,\dd t.
\label{eq:burstintegral}
\end{equation}
A direct discretisation makes this plain: partitioning time into intervals
of width $\Delta t$, the probability that the burst occurs during an interval
in which the process is at $k$ is
$$
\pr{\xt=k}\,\pr{\text{rupture in }\Delta t\mid\xt=k}
=\pr{\xt=k}\,\delta k\,\Delta t,
$$
and summing over intervals and passing to the limit $\Delta t\to0$ yields
\eqref{eq:burstintegral}.

The integral is evaluable in closed form. With $P(t)$ as in \eqref{eq:Pdef},
\begin{equation}
\frac{\dd}{\dd t}\Bigl[\frac{P(t)^k}{k\lambda}\Bigr]
=\frac{1}{\lambda}P(t)^{k-1}P'(t)
=\lrb{\frac{a-b}{b-a\ee^{\lambda(a-b)t}}}^2
\ee^{\lambda(a-b)t}P(t)^{k-1}
=p_k(t),
\label{eq:antideriv}
\end{equation}
as one verifies from
$P'(t)=\lambda(a-b)^2\ee^{\lambda(a-b)t}/(b-a\ee^{\lambda(a-b)t})^2$.
Therefore, since $P(0)=0$ and $P(\infty)=1/a$,
\begin{equation}
\int_0^\infty p_k(t)\,\dd t=\frac{1}{k\lambda}\Bigl(\frac{1}{a}\Bigr)^k,
\end{equation}
and \eqref{eq:burstintegral} gives the burst-size law.

\begin{theorem}[Burst-size distribution]
\label{thm:burst}
Unconditionally (including the mass $b$ of realisations that never burst),
\begin{equation}
\pr{\Kb=k}=\frac{\delta}{\lambda}\,a^{-k},\qquad k\ge1;
\label{eq:burstlaw}
\end{equation}
conditioned on the burst occurring,
\begin{equation}
\pr{\Kb=k\mid\text{burst}}=(a-1)\,a^{-k},\qquad k\ge1.
\label{eq:burstlaw-cond}
\end{equation}
\end{theorem}

\begin{proof}
The unconditional form is the computation above. For the conditional form,
divide \eqref{eq:burstlaw} by the eventual-burst probability $1-b$:
$$
\pr{\Kb=k\mid\text{burst}}
=\frac{(\delta/\lambda)\,a^{-k}}{1-b}
=\frac{AB}{B}\,a^{-k}
=(a-1)\,a^{-k},
$$
using $\delta/\lambda=AB$ from \eqref{eq:AB}.
\end{proof}

No simulation is required; nevertheless two consistency checks close the
loop pleasantly. First, normalisation:
$$
\sum_{k=1}^{\infty}\frac{\delta}{\lambda}a^{-k}
=\frac{\delta}{\lambda}\,\frac{1}{a-1}
=\frac{AB}{A}=B=1-b,
$$
matching the eventual-burst probability. Second, the first moment:
$$
\sum_{k=1}^{\infty}k\,\frac{\delta}{\lambda}a^{-k}
=\frac{\delta}{\lambda}\,\frac{a}{(a-1)^2}
=\frac{aB}{A}=V_\infty,
$$
recovering \eqref{eq:Vinf}. This second check independently validates the
burst-size law, the formula for $V_\infty$, and the identity
\eqref{eq:AB} simultaneously.

\subsection{Burst size equals the quasi-stationary distribution}
\label{sec:burst=qsd}

Compare \eqref{eq:burstlaw-cond} with Theorem~\ref{thm:qsd}: they are the
same law.

\begin{corollary}[Burst size $=$ QSD]
\label{cor:burst-qsd}
The burst-size distribution, conditioned on bursting, is exactly the
quasi-stationary distribution of the load: both are geometric on
$\{1,2,\ldots\}$ with ratio $1/a$.
\end{corollary}

This is not obvious, and it deserves a moment's pause. A burst is a
\emph{size-biased} draw from the transient load distribution, taken against
the burst intensity, and then integrated over all burst times; that this
should land exactly on the conditional limiting distribution is a genuine
identity, not a definition. The algebra above proves it; a structural
argument --- presumably relating size-biasing composed with the burst
intensity to the Yaglom limit \cite{yaglom1947certain} --- would make it
quotable rather than coincidental, and remains open
(see \S\ref{sec:open-problems}).

\subsection{Size-biasing and the size of a late burst}
\label{sec:size-bias}

The size-biasing mechanism is worth stating on its own terms. A cell bursts
at rate $\delta\xt$, so given that the burst occurs at time $t$, its size is
a size-biased draw from $p_\cdot(t)$:
\begin{equation}
\pr{\Kb=k\mid\tau=t}=\frac{k\,p_k(t)}{J(t)},
\qquad
\mean{\Kb\mid\tau=t}=\frac{K(t)}{J(t)}
=1+\frac{2\lambda}{\delta}\lrb{1-I(t)}.
\label{eq:sizebias}
\end{equation}
Taking $t\to\infty$,
\begin{equation}
\mean{\Kb\mid\tau=t}\xrightarrow[t\to\infty]{}
1+\frac{2\lambda}{\delta}(1-b)
=\frac{\langle X^2\rangle_{\rm QS}}{\langle X\rangle_{\rm QS}},
\end{equation}
where the second equality follows from \eqref{eq:QSmoments}. A very late
burst is roughly twice the size of an average one --- the signature of
size-biasing a geometric law: the cells that survive long are precisely the
ones that grew large, and large cells burst quickly.

\figureflag{F4a.5}{Burst size and late bursts}{%
\textbf{Purpose:} show the two identities of this section together --- the
burst-size law, and the size-biasing effect by which late bursts exceed
average ones. \textbf{Type:} Python analytic plot, two panels.
\textbf{Content:} panel (a): bars of the unconditional burst-size law
$\Pr\{\Kb=k\}=(\delta/\lambda)a^{-k}$, $k=1,\ldots,40$, for
$(\lambda,\mu,\delta)=(1,0.2,0.05)$, with the conditional law
$(a-1)a^{-k}$ overlaid as a connected curve, and the mass $b$ at $k=0$
annotated (bar or text). Panel (b): the size-biased conditional mean
$K(t)/J(t)=1+\frac{2\lambda}{\delta}(1-I(t))$ against $t\in[0,25]$ for the
same parameters, with the horizontal asymptote
$(a+1)/(a-1)\approx33.5$ dashed and annotated; optionally overlay the QS
mean $a/(a-1)$ for contrast. \textbf{Axes:} panel (a): burst size $k$
(horizontal), probability (vertical); panel (b): $t$ (horizontal),
$\mean{\Kb\mid\tau=t}$ (vertical). \textbf{Style:} unified Python style ---
default matplotlib with \texttt{tab:} colours, grid alpha 0.25,
\texttt{tight\_layout}, PDF output, fonts matching the chapter.
\textbf{Data source:} analytic formulas \eqref{eq:burstlaw},
\eqref{eq:burstlaw-cond}, \eqref{eq:sizebias}; no simulation.
\textbf{Placement:} \Secref{sec:size-bias}, after the size-biasing
discussion.}

\subsection{Burst time given burst size, without death}
\label{sec:tau-given-k}

Size-biasing conditions the burst size on the burst time; the converse
conditioning --- the burst time given the burst size --- has an equally
explicit form when $\mu=0$, and it is the quantity an experiment accesses
when it measures both members of each (burst time, burst size) pair. Write
$\theta=\lambda+\delta$ for the rate combination that governs the $\mu=0$
case (it is $\lambda(a-b)$ evaluated at $b=0$). Fixing $n$ and regarding
$p_n(t)$ as a function of $t$, the state probability rises from zero,
peaks, and decays; differentiating $\log p_n(t)$ shows the maximum is
attained at
\begin{equation}
t_n=\frac{\log n}{\theta}:
\label{eq:tn}
\end{equation}
a burst of size $n$ is most likely to have occurred once the process has had
of order $(\log n)/\theta$ of growth time. The exact conditional mean burst
time is the following.

\begin{proposition}[Conditional rupture time, $\mu=0$]
\label{prop:tau-given-k}
With $\mu=0$ and $\theta=\lambda+\delta$,
\begin{equation}
\mean{\tau\mid\Kb=n}=\frac{1}{\theta}\sum_{k=1}^{n}\frac{1}{k},
\qquad n\ge1.
\label{eq:tau-given-k}
\end{equation}
\end{proposition}

\begin{proof}
When $\mu=0$ the state probabilities \eqref{stateprob} reduce, with
$r=\lambda/\delta$ and $s=\theta t$, to
$$p_n(t)=\ee^{-\theta t}\lrb{\frac{1-\ee^{-\theta t}}{1+1/r}}^{n-1},$$
and the burst-size law \eqref{eq:burstlaw} to
$\pr{\Kb=n}=(1/r)(r/(r+1))^n$. Writing
$$L_1(n)=\int_0^\infty \ee^{-s}(1-\ee^{-s})^{n-1}\,\dd s,\qquad
L_2(n)=\int_0^\infty s\,\ee^{-s}(1-\ee^{-s})^{n-1}\,\dd s,$$
the substitution $s=\theta t$ gives
$\mean{\tau\mid\Kb=n}=L_2(n)/(\theta L_1(n))$. The first integral is
$L_1(n)=1/n$, by the substitution $u=1-\ee^{-s}$. For the second, split off
one factor of $(1-\ee^{-s})$ to write
$L_2(n+1)=L_2(n)-\int_0^\infty s\ee^{-2s}(1-\ee^{-s})^{n-1}\dd s$, and
integrate by parts in the remaining integral; the result is the recurrence
$$(n+1)L_2(n+1)=nL_2(n)+\frac{1}{n+1},\qquad L_2(1)=1,$$
whose solution is $L_2(n)=\frac{1}{n}(1+\frac12+\cdots+\frac1n)$; the
quotient is \eqref{eq:tau-given-k}.
\end{proof}

For large $n$ the harmonic sum gives the approximation
\begin{equation}
\mean{\tau\mid\Kb=n}\simeq\frac{\log n+\gamma}{\theta}\qquad(n\gg1),
\label{eq:tau-given-k-asymp}
\end{equation}
where $\gamma=0.5772\ldots$ is Euler's constant: conditional rupture time
grows logarithmically with burst size. The full conditional distribution has
a clean form in the same limit. Since the density of $\tau$ at $t$, given
$\Kb=n$, is proportional to $p_n(t)$,
\begin{equation}
\pr{\tau<t\mid\Kb=n}
=\frac{\delta n\int_0^t p_n(s)\,\dd s}{\pr{\Kb=n}}
=(1-\ee^{-\theta t})^n,
\label{eq:cond-cdf}
\end{equation}
and for $n\gg1$, with $T=\theta(t-t_n)$ and $t_n$ as in \eqref{eq:tn},
\begin{equation}
\pr{\tau<t\mid\Kb=n}\simeq\exp\lrb{-n\ee^{-\theta t}}
=\exp\lrb{-\ee^{-T}}:
\end{equation}
the conditioned rupture time, centred at $t_n$ and scaled by $1/\theta$,
converges to the standard Gumbel law. Burst sizes thus select burst times
from a travelling window of width of order $1/\theta$ around
$(\log n)/\theta$.

\subsection{The mean burst time when death is possible}
\label{sec:tau-burst-mu-pos}

When $\mu>0$ the burst time $\tau$ has no finite mean, because the mass $b$
at $\tau=\infty$ (the internally cleared cells) makes the integral diverge.
Conditioned on bursting, however, the mean is finite and explicit.

\begin{proposition}[Mean burst time conditioned on bursting]
\label{prop:tau-burst}
For $\mu>0$,
\begin{equation}
\mean{\tau\mid\text{burst}}
=\frac{1}{\lambda(1-b)}\log\frac{a-b}{a-1}.
\label{eq:tau-burst}
\end{equation}
\end{proposition}

\begin{proof}
The density of $\tau$ on the bursting realisations is $\varphi(t)=\delta J(t)$,
with total mass $1-b$, so
$\mean{\tau\mid\text{burst}}=(1-b)^{-1}\int_0^\infty t\,\varphi(t)\,\dd t$.
Since $\varphi=-(I-b)'$, integrating by parts on $[0,T]$ gives
$$\int_0^T t\,\varphi(t)\,\dd t
=-T\bigl(I(T)-b\bigr)+\int_0^T\bigl(I(t)-b\bigr)\,\dd t,$$
and the boundary term vanishes as $T\to\infty$ because
$I(T)-b=B(a-b)/(B+Aw)$, read off from \eqref{eq:IDIhat}, decays like
$\ee^{-\theta T}$. It remains to evaluate the integral: with
$\dd t=\dd w/(\theta w)$ and the partial-fraction decomposition
$1/[w(B+Aw)]=B^{-1}\bigl(1/w-A/(B+Aw)\bigr)$,
$$\int_0^\infty\bigl(I(t)-b\bigr)\,\dd t
=\frac{B(a-b)}{\theta}\int_1^\infty\frac{\dd w}{w(B+Aw)}
=\frac{1}{\lambda}\log\frac{a-b}{a-1},$$
and dividing by $1-b$ gives \eqref{eq:tau-burst}.
\end{proof}

This quantity is distinct from the mean productive lifetime
$\mean{T_{\rm prod}}=\lambda^{-1}\log\frac{a}{a-1}$ of
\S\ref{sec:lifetime}, and the distinction should be kept visible:
$\mean{\tau\mid\text{burst}}$ averages the rupture time over cells that
burst, while $\mean{T_{\rm prod}}$ averages the time spent productive over
\textit{all} cells, the internally cleared ones included with their shorter
productive spans. When $\mu=0$ the two coincide, both reducing to
$\lambda^{-1}\log(1+\lambda/\delta)$, and it is tempting to read more into
the coincidence than it supports.

\subsection{Comparison with budding release}
\label{sec:bud-vs-burst-single-cell}

The bursting single-cell picture just completed has a natural comparator:
the budding cell of the constant-rate models. In its simplest stochastic
form, a budding infected cell experiences two \textit{independent} event
types --- release of one infectious particle at rate $p$, and death of the
cell at rate $d_{\Icell}$ --- and the two continue until the cell dies. Its
survival and mean release are then
$$I_{\rm bud}(t)=\ee^{-d_{\Icell}t},\qquad
V_{\rm bud}(t)=\frac{p}{d_{\Icell}}\lrb{1-\ee^{-d_{\Icell}t}},$$
the same one-cell limit that the Basic Model of Viral Replication reaches at
the population level. The release count of a budding cell is geometric on
$\{0,1,2,\ldots\}$:
\begin{equation}
\pr{\Kb_{\rm bud}=k}=\frac{d_{\Icell}}{d_{\Icell}+p}
\lrb{\frac{p}{d_{\Icell}+p}}^k,\qquad k=0,1,2,\ldots,
\label{eq:bud-geom}
\end{equation}
since each of the successive events is a release with probability
$p/(d_{\Icell}+p)$ and the terminating death with probability
$d_{\Icell}/(d_{\Icell}+p)$. The structural contrast with bursting is exact:
in budding, release and death are independent events, and a cell may release
many particles and then die with nothing left to release; in bursting,
release and death are the \textit{same} event, and whatever has accumulated
is released at once. Table~\ref{tab:budburst} collects the two single-cell
descriptions side by side for the $\mu=0$ bursting case.

\begin{table}[H]
\centering
\renewcommand{\arraystretch}{1.4}
\small
\begin{tabular}{@{}lll@{}}
\toprule
 & budding & bursting ($\mu=0$) \\
\midrule
$\ddt{I}$ & $-d_{\Icell} I$ & $-(\lambda+\delta)I+\lambda I^2$ \\
$I(t)$ & $\ee^{-d_{\Icell}t}$ & $\dfrac{\lambda+\delta}{\lambda+\delta\,\ee^{(\lambda+\delta)t}}$ \\
mean cell lifetime & $1/d_{\Icell}$ & $\dfrac{1}{\lambda}\log\lrb{1+\dfrac{\lambda}{\delta}}$ \\
$V(t)$ & $\dfrac{p}{d_{\Icell}}(1-I)$ & $\dfrac{\lambda}{\delta}+1-\dfrac{2\lambda+\delta}{\delta}I+\dfrac{\lambda}{\delta}I^2$ \\
release law & geometric on $\{0,1,\ldots\}$, \eqref{eq:bud-geom} & geometric on $\{1,2,\ldots\}$, $\pr{\Kb=k}=\frac{\delta}{\lambda}a^{-k}$ \\
\bottomrule
\end{tabular}
\caption{Single-cell budding and bursting compared, $\mu=0$ for bursting.
Budding parameters: release rate $p$ and infected-cell death rate
$d_{\Icell}$. Bursting parameters: division rate $\lambda$ and rupture rate
$\delta$ per intracellular particle; $a=(\lambda+\delta)/\lambda$ here.}
\label{tab:budburst}
\end{table}

Two rows of the table carry the comparison. The mean cell lifetime of a
bursting cell grows only logarithmically with $\lambda/\delta$, while its
mean release grows linearly; budding offers no such coupling, its lifetime
and its release controlled by separate parameters. And the release laws,
both geometric, differ in their support: budding can release nothing at all
from a dead cell, while bursting never does. These two differences are what
the population-level comparison of Chapter~4b exploits, where bursting's
simultaneous release-and-death translates into a measurable advantage at
small population sizes.

\figureflag{F4a.8}{Conditional rupture time versus burst size}{%
\textbf{Purpose:} show the harmonic-number law \eqref{eq:tau-given-k}
against burst size, with its large-$n$ approximation and the linear budding
comparator, in one figure. \textbf{Type:} Python analytic plot.
\textbf{Content:} main panel: $\mu=0$ with $\theta=\lambda+\delta=1.1$
(e.g.\ $\lambda=1$, $\delta=0.1$): red dots at the exact means
$\mean{\tau\mid\Kb=n}=\theta^{-1}\sum_{k=1}^n 1/k$ for $n=1,\ldots,40$;
dotted curve $\theta^{-1}(\log n+\gamma)$ for $n\ge2$; dashed straight line
$(n+1)/(d_{\Icell}+p)$, the budding conditional mean for exactly $n$
released particles, with $d_{\Icell}+p=8$ (the comparator scale of the
source figure; any matched scale may be substituted and should then be
stated in the caption). Optional second panel: three values of $\mu$ ---
$\lambda=0.5$, $\delta=0.01$, $\mu\in\{0,0.005,0.01\}$ --- exact conditional
means from simulation of the conditioned process, with the approximation
$(\log n+\gamma-\mu/\lambda)/(\lambda+\delta-\mu)$ overlaid as a dashed
curve for each. \textbf{Axes:} burst size $n$ (horizontal), conditional mean
rupture time (vertical); legend entries ``exact'', ``harmonic
approximation'', ``budding''. \textbf{Style:} unified Python style ---
default matplotlib with \texttt{tab:} colours, grid alpha 0.25,
\texttt{tight\_layout}, PDF output, fonts matching the chapter.
\textbf{Data source:} exact formula \eqref{eq:tau-given-k}; approximation
\eqref{eq:tau-given-k-asymp}; budding comparator: $n+1$ exponential waiting
times each of rate $d_{\Icell}+p$; optional panel by simulating the BDC
process and conditioning on $\Kb=n$. \textbf{Placement:}
\Secref{sec:tau-given-k}, after Proposition~\ref{prop:tau-given-k}.}
